\documentclass[a4paper,10pt]{article}
\usepackage[utf8]{inputenc}
\usepackage{amssymb}
\usepackage{amsmath}
\usepackage{graphicx}
\newcommand{\HRule}{\rule{\linewidth}{0.5mm}}
\usepackage{subfigure}
\usepackage{multicol}
\usepackage[usenames,dvipsnames]{color}
\definecolor{darkgray}{rgb}{0.95,0.95,0.95}
\usepackage{listings}
\usepackage{color}
\usepackage{textcomp}
\definecolor{listinggray}{gray}{0.9}
\definecolor{lbcolor}{rgb}{0.9,0.9,0.9}
\lstset{
    backgroundcolor=\color{lbcolor},
    tabsize=4,
    rulecolor=,
    language=C++,
    basicstyle=\scriptsize,
    upquote=true,
    aboveskip={1.5\baselineskip},
    columns=fixed,
    showstringspaces=false,
    extendedchars=true,
    breaklines=true,
    prebreak = \raisebox{0ex}[0ex][0ex]{\ensuremath{\hookleftarrow}},
    frame=single,
    showtabs=false,
    showspaces=false,
    showstringspaces=false,
    identifierstyle=\ttfamily,
    keywordstyle=\color[rgb]{0,0,1},
    commentstyle=\color[rgb]{0.133,0.545,0.133},
    stringstyle=\color[rgb]{0.627,0.126,0.941},
    backgroundcolor=\color{darkgray},
}

\usepackage[left=1.0cm, right=1.0cm, top=2cm, bottom=4cm]{geometry}
\usepackage{fancyhdr}
\pagestyle{fancy}
\usepackage{lastpage}
\renewcommand\headrulewidth{1pt}
\fancyhead[L]{\textsc{Nachos Étape 3 : Multithreading}}
\fancyhead[R]{\textsc{Polytech' Grenoble}}
\renewcommand\footrulewidth{1pt}
\fancyfoot[R]{ \textsc{RICM 4}}

\begin{document}
\begin{titlepage}

\begin{center}


% Upper part of the page
\includegraphics[width=0.25\textwidth]{./images/logo}\\[1cm]

\textsc{\LARGE Polytech' Grenoble}\\[1.5cm]

\textsc{\Large RICM 4\`eme ann\'ee}\\[1.2cm]


% Title
\HRule \\[0.4cm]
{ \huge \bfseries NachOS\\[0.6cm]
Etape 3: Multithreading}
\\[0.4cm]

\HRule \\[2cm]

% Author and supervisor
\begin{minipage}{0.4\textwidth}
\begin{flushleft} \large
\emph{\'Etudiants:}\\
Elizabeth \textsc{Paz} \\
Salem \textsc{Harrache}
\end{flushleft}
\end{minipage}
\begin{minipage}{0.4\textwidth}
\begin{flushright} \large
\emph{Professeur:} \\
Vania \textsc{Marangozova}
\end{flushright}
\end{minipage}

\vfill

% Bottom of the page
{\large  février 2012}

\end{center}

\end{titlepage}

\section{Git}

Pour voir les differences entre l'etape 2 et l'etape 3 vous pouvez lancer un
diff avec le tag step2 :
\begin{lstlisting}
git diff step2
\end{lstlisting}
ou alors directement avec le commit 4a23b1..
\begin{lstlisting}
git diff 4a23b18e1b2b597b0454b768347d464f47ecb572
\end{lstlisting}

\section{Test de l'étape 3}

Vous pouvez soit lancer le script de test automatique \textit{runtest.sh} soit
lancer manuellement les deux programmes \textit{./build-origin/nachos-userprog
-rs 1 -x ./build/jointhreads} et \textit{./build-origin/nachos-userprog -rs 1 -x
./build/makethreads}

\begin{lstlisting}
$ ./runtest.sh
Compilation en cours...
Lancement du Test 1 : Create & Exit :)
--------------------------------------
./build-origin/nachos-userprog -rs 1 -x ./build/makethreads
Test de thread :
Thread : 0 OK
Thread : 1 OK
Thread : 2 OK
Thread : 3 OK
Thread : 4 OK
Thread : 5 OK
Thread : 6 OK
Thread : 7 OK
Thread : 8 OK
Impossible de créer de nouveaux Thread
Fin du main :
0
2
1
3
4
6
5
8
7
Machine halting!
Fin Test 1
Lancement du Test 2 : Join :)
-----------------------------
./build-origin/nachos-userprog -rs 1 -x ./build/jointhreads
Je lance le Thread 1
1
Fin Thread 1
Je lance le Thread 2
2
Fin Thread 2
Je lance le Thread 3
3
Fin du main :
Machine halting!
\end{lstlisting}

Sans l'option -rs fait le thread  main ne donne jamais la main aux autres
threads. Comme on peut le voir sur le test 1, le nombre de thread simultané
est limité (par la taille de l'espace d'adressage et la taille de la
pile d'un thread). Dans ce cas un code d'erreur (-1) est renvoyé lors de
l'appel à \textit{UserThreadCreate}.

\section{Choix de conception}

\subsection{La classe UserThread}
Pour implementer les threads utilisateurs nous avons créé une nouvelle classe
UserThread qui hérite de la classe Thread. Cette classe surcharge le
comportement de la methode Thread::Fork(..) et implemente quelques mécanismes
utiles pour les threads utilisateurs. Cet objet va servir également à stocker
les arguments qu'on passe à la fonction UserThreadCreate(), ce qui nous evite
de devoir utiliser une structure dédié au passage de parametre à Fork.

\textit{code/userprog/userthread.h}
\begin{lstlisting}
class UserThread : public Thread {
    public:
        UserThread(const char *name, int f, int a, int callback);
        int func;
        int arg;
        void Fork (); // Make userthread run (*f)(arg)
        void UpdateCallBackRegister (int value); // $31 = value
};
\end{lstlisting}

\textit{code/userprog/userthread.cc}
\begin{lstlisting}
void StartUserThread(int thread) {
    UserThread *t = (UserThread *) thread;
    // L'id du thread informe egalement le numéro de page du thread
    currentThread->space->InitThreadRegisters(t->func, t->arg, t->getId());
    currentThread->space->UpdateRunningThreads(1); // appel atomique
    machine->Run();
}

void UserThread::Fork () {
    DEBUG ('t', "Forking userThread \"%s\"\n", getName ());
    Thread::Fork (StartUserThread, (int) this);
}
\end{lstlisting}

\subsection{Appel implicite de UserThreadExit}

La fonction UpdateCallBackRegister permet de modifier l'adresse de retour
(registre 31 dans MIPS) lors de l'appel système. En effet pour quitter les
threads sans appel explicite de la fonction UserThreadExit(), nous devons au
moment du lancement du nouveau thread de lui affecter comme adresse de retour
l'adresse de la fonction UserThreadExit qu'on passe en paramètre dans
le fichier start.S (registre 6)

\begin{lstlisting}
    .globl UserThreadCreate
    .ent UserThreadCreate
UserThreadCreate:
    addiu $2,$0,SC_UserThreadCreate
    addiu $6,$0,UserThreadExit // On passe en parametre la fonction de retour
    syscall
    j    $31
    .end UserThreadCreate
\end{lstlisting}


\subsection{Modifications de l'espace d'adressage du processus}

Pour gérer plusieurs threads, nous sommes obligé de modifier l'espace
d'adressage (AddrSpace) du processus. Nous avons entre autre ajouté le nombre
de thread en execution, un objet bitmap pour gérer l'allocation des
zones des threads et des sémaphores pour utiliser ces attribus en section
critique. Chaque thread a sa propre pile qui est une partie de la pile du
processus (Sa taille est de 3 pages). Nous somme donc naturellement limité en
nombre de threads en excution simultanée par processus. Lorsque cette limitation
se manifeste, on renvoie simplement un code d'erreur (-1).

Nous devons également differencier un thread d'un autre, pour celà on ajoute un
attribut \textbf{Thread::id} et un compteur de thread pour eviter de réutiliser
les identifiants de threads. On créer ensuite une association entre un thread
(son id) et le numéro de la zone auquel il sera affecté lors de la création de
ce dernier. Bien evidement la gestion de ces deux map (\textbf{bitmap} et
\textbf{id//zone}) se fait en section critique.

L'espace d'adressage est la partie commune de tous les threads (d'un même
processus). C'est pour cette raison que nous avons ajouté plusieurs methodes
qui vont ``rendre service'' aux threads. Par exemple, pour savoir s'il est le
seul en cours d'execution, un thread ne peut pas le savoir directement, il va
utiliser la méthode \textbf{int AddrSpace::Alone()} qui va garantir que cette
appel est atomique. C'est à AddrSpace que revient la responsabilité de garder
un état cohérent des variable partagées.

\begin{lstlisting}
int AddrSpace::Alone() {
    int value = 0;
    this->semRunningThreads->P();
    if (this->runningThreads == 0)
        value = 1;
    this->semRunningThreads->V();
    return value;
}
\end{lstlisting}

\subsection{Le Thread main est comme tous les autres}

Notre implémentation des threads utilisateurs considere que même le thread main
se comporte comme un thread utilisateur, c'est à dire qu'il sa propre pile (3
pages) et qu'il cède sa place (sa pile) aux autres threads une fois terminé. Ce
qui as pour principale consequence que l'appel \textbf{Exit()} ne se fait plus
automatiquement à la fin du main, mais sera fait par le dernier thread en cours.

Ceci dit, un thread (main ou pas) peut toujours appeller explicitement
\textbf{Exit()} pour arrêter le processus.


\subsection{Mécanisme d'attente des threads (UserThreadJoin)}

Nous utilisons un tableau de sémphore de la taille du nombre maximal de threads.
Chaque thread (A) qui veut se bloque sur un autre (B) va devoir prendre un jeton
du sémaphore lié au thread B sur lequel il veut se bloquer. Ces sémaphore sont
donc bloquant lorsque le thread B est en cours d'execution.

Une fois terminé, le thread B va reveiller les threads qui se sont bloqué que
lui en deposant un jeton dans le sémphore concerné.
Pour permettre à plusieurs thread de se bloquer sur le meme thread , la phase
de reveil doit se faire en chaîne, c'est à dire chaque thread reveillé doit
reveillé le suivant etc..

Ce mécanisme doit également bloqué le prochain thread qui va s'exécuter dans
la nouvelle zone liberé par le thread B.

Pour arriver à se resultat, nous initialisons le tableau de sémaphore à 1 jeton
chacun. A chaque creation, il faut prendre un jeton (pour le sémphore devienne
bloquant pour tout les threads qui appelle \textbf{UserThreadJoin()}). Si il y
a une phase de reveil qui est en cours, ce nouveau thread sera bloqué, et
c'est ce qu'on cherche pour pas qu'un thread bloqué sur un thread spécifique
se retrouve bloqué sur d'autres (avec le risque qu'il ne reveille jamais)

\begin{lstlisting}
int do_UserThreadCreate(int f, int arg, int callback) {
    [...]
    // Avant de commencer on prend le jetton, pour que tout thread qui appelle
    // userThreadJoin sur moi soit bloqué.
    currentThread->space->semJoinThreads[newThread->getZone()]->P();
    [...]
    newThread->Fork();
    return newThread->getId();
}
\end{lstlisting}

\begin{lstlisting}
void do_UserThreadExit() {
    currentThread->space->UpdateRunningThreads(-1); // appel atomique
    [...]
    // Je libere les threads en attente sur moi
    currentThread->space->semJoinThreads[currentThread->getZone()]->V();
    // Plusieurs threads peuvent attendre que je me termine.
    // Il faut donc que dans la fonction join, les threads en attente se
    // reveillent les uns les autres
    [...]
    currentThread->Finish();
}
\end{lstlisting}

\begin{lstlisting}
int do_UserThreadJoin(int thread_id) {
    int zone = currentThread->space->GetZoneFromThreadId(thread_id);
    if (zone < 0)
        return zone;
    currentThread->space->semJoinThreads[zone]->P();
    // On reveille le suivant qui peut etre soit le prochain thread à qui
    // on a allouer la zone, soit un autre thread qui avait appeller join
    currentThread->space->semJoinThreads[zone]->V();
    return 0;
}
\end{lstlisting}


\subsection{Améliorations possibles}

\subsubsection{Les identifiants de threads}

Pour éviter de réutiliser les identifiants de threads, nous avons mis en place
un compteur qu'on incremente. Il faudrai prévoir un mécanisme qui remet à zéro
ce compteur.

\subsubsection{Gestion de l'espace des threads}

Un handicap important de notre implmentation est la gestion des dépassements
mémoire : on ne garantis à \textbf{aucun moment} qu'un thread ne va pas depasser
sa pile, et donc il peut tres bien écraser la mémoire d'un autre thread d'un
autre thread. Il faudrais mettre en place des mecanisme pour limiter
l'utilisation de la pile d'un thread a son espace propre. Par exemple, vérfier a
chaque lecture memoire, qu'il est autoriser à ecrire/lire.

Une des amélorations possibles, serait la gestion dynamique de la taille
de la pile d'un thread. Dans notre solution (3 pages par thread) On a soit de
la fragmentation interne soit de depassements, mais meme avec un mecanisme de
non-depassement, le thread peut planter, car il n'a pas la possibilité
d'etendre sa pile. Une solution possible est d'utiliser pour
chaque nouveau thread une taille de pile minimale (une page), mais en cas de
depassement, ajouter une nouvelle page a sa pile...Ce qui revient à prévoir un
mécanisme de segmentation inter-processus

\end{document}

